% part: intuitionistic-logic
% chapter: propositions-as-types
% section: proofs-to-terms

\documentclass[../../../include/open-logic-section]{subfiles}

\begin{document}

\olfileid[tr]{int}{pty}{pt}

\olsection{\usetoken{P}{derivation}{}den İspat Terimlerine Geçiş}

Doğal tümdengelim !!{proof}s{}ını ispat terimlerine dönüştürme, yani
\Log{N2i}'deki !!{proof}s{}ı \Log{tN2i}'deki !!{proof}s{}a çevirme süreci son
derece doğrudandır. Yalnızca !!{derivation} içindeki her sekantın sağındaki
!!{formula}{}ün soluna bir ispat terimi yazılır; böylece $\Gamma \Sequent M
: !A$ biçiminde sekantlar elde edilir. Bundan sonra $M$'nin $\Gamma$
bağlamında~$!A$'ya \emph{tanıklık ettiğini} söyleyeceğiz. Hangi ispat
teriminin yazılacağı bütünüyle \Log{tN2i}'nin kurallarıyla belirlenir.

\begin{prop}
\Log{N2i}'de $\Gamma \Sequent !A$'nın her !!{proof}~$\delta$'sına,
\Log{tN2i}'de $\Gamma \Sequent N : !A$'nın bir !!a{proof}~$\delta'$ karşılık
gelir. İspat terimi $N$,~$\delta$ tarafından biricik biçimde belirlenir.
\end{prop}

\begin{proof}
$\delta$'nın yüksekliği üzerinde tümevarımla ilerleriz. $\delta$, $x:!A
\Sequent !A$ aksiyomuysa $\delta'$,~$x:!A \Sequent x:!A$ aksiyomudur.
$\delta$ bir çıkarım kuralıyla bitiyorsa tümevarım varsayımı her öncül için
ispat terimleri ve her öncüle karşılık gelen bu ispat terimlerini içeren
sekantların \Log{tN2i}-!!{proof}s{}ını verir. \Log{tN2i}'nin karşılık gelen
kuralını uygularız; karşılık gelen ispat terimi \Log{tN2i} çıkarım kuralıyla
belirlenir.
\end{proof}

\begin{ex}
  \Log{N1i}'de $(!A \land !B) \lif (!B \land !A)$'nın şu çok basit
  !!{proof}{}ını ele alalım:
  \[
  \AxiomC{$\Discharge{!A\land !B}{x}$}
  \RightLabel{\Elim\land[2]}
  \UnaryInfC{$!B$}
  \AxiomC{$\Discharge{!A\land !B}{x}$}
  \RightLabel{\Elim\land[1]}
  \UnaryInfC{$!A$}
  \RightLabel{\Intro\land}
  \BinaryInfC{$!B \land !A$}
  \DischargeRule{\Intro\lif}{x}
  \UnaryInfC{$(!A \land !B) \lif (!B \land !A)$}
  \DisplayProof
  \]
  \Log{N2i}'de şöyle görünür:
  \[
  \Axiom$x : !A\land !B \fCenter !A\land !B$
  \RightLabel{\Elim\land[2]}
  \UnaryInf$x : !A\land !B \fCenter !B$
  \Axiom$x : !A\land !B \fCenter !A\land !B$
  \RightLabel{\Elim\land[1]}
  \UnaryInf$x : !A\land !B \fCenter !A$
  \RightLabel{\Intro\land}
  \BinaryInf$x : !A\land !B \fCenter !B \land !A$
  \RightLabel{\Intro\lif}
  \UnaryInf$\fCenter (!A \land !B) \lif (!B \land !A)$
  \DisplayProof
  \]
  İspat terimlerini yukarıdan aşağıya atarız. Aksiyomların ispat terimleri,
  bağlamdaki değişkenin kendisi olan~$x$'tir. $\Elim\land[i]$ ile ilişkili
  ispat terimleri böylece $\proj{2}{x}$ ve~$\proj{1}{x}$ olarak belirlenir.
  \Intro{\land} uygulandığında bu iki ispat terimini birleştiririz:
  $\pair{\proj{2}{x}}{\proj{1}{x}}$. Son olarak \Intro{\lif} kuralı bir
  $\lambd$-soyutlaması uygulayıp~$x$'i bağlar ve $x$'in $!A \land !B$
  varsayımını etiketlediğini kaydeder:
  $\lambd[\typeof{x}{!A \land !B}][\pair{\proj{2}{x}}{\proj{1}{x}}]$.
  O zaman \Log{tN2i}'deki !!{proof} şöyledir:
  \[
  \Axiom$x : !A\land !B \fCenter x: !A\land !B$
  \RightLabel{\Elim\land[2]}
  \UnaryInf$x : !A\land !B \fCenter \proj{2}{x} : !B$
  \Axiom$x : !A\land !B \fCenter x: !A\land !B$
  \RightLabel{\Elim\land[1]}
  \UnaryInf$x : !A\land !B \fCenter \proj{1}{x} : !A$
  \RightLabel{\Intro\land}
  \BinaryInf$x : !A\land !B \fCenter \pair{\proj{2}{x}}{\proj{1}{x}} : 
    !B \land !A$
  \RightLabel{\Intro\lif}
  \UnaryInf$\fCenter \lambd[\typeof{x}{!A \land
  !B}][\pair{\proj{2}{x}}{\proj{1}{x}}] : (!A \land !B) \lif (!B \land !A)$
  \DisplayProof
  \]
\end{ex}


\end{document}
